\chapter{Contexte général du projet}
\clearpage

\section*{Introduction}
\addcontentsline{toc}{section}{Introduction}

À travers ce chapitre, nous présenterons l'environnement dans lequel ce stage s'est déroulé, la problématique que nous avons traitée et la méthodologie de travail suivie.

\section{Organisme d’accueil}

\subsection{Présentation générale de SQLI DIGITAL EXPERIENCE}

\textbf{SQLI Digital Experience} est une Société de Services et d’Ingénierie en Informatique, multinationale, spécialisée dans l'accompagnement des entreprises dans l'utilisation des nouvelles technologies. 

\begin{figure}[!h]
    \centering
    \includegraphics[width=6cm]{images/logo.png}
    \caption{Logo de SQLI Digital Experience \cite{sqli}}
    \label{fig:my_label}
\end{figure}

SQLI est créée en 1990 par Alain Lefebvre et Jean Rouveyrol. Dès lors, elle aide les entreprises à adopter les technologies disruptives et les méthodologies modernes, à créer des expériences client omnicanales avec des solutions web et mobiles, des plateformes e-commerce, des solutions de marketing digital et des solutions Customer Relationship Management\cite{sqli}. 

En 2002, SQLI a été la première société de services européenne à intégrer le modèle de qualité "Capability Maturity Model Integration"\cite{cmmi} pour sa production. 
 
\subsection{Organisation interne}

Les agences de SQLI au Maroc forment l'entité Innovative Services Center Maroc, sous la direction d'Eric Chanal. 

\subsection{Secteurs d'activités}

Ses pôles métiers sont : 
\begin{itemize}
    \item Le pôle « stratégie en systèmes d’information »
    \item Le pôle « conception web : studio SQLI »
    \item Le pôle « formation et transfert de compétences »
    \item Le pôle « ingénierie et intégration »
\end{itemize}

\subsection{Valeurs de SQLI}


\section{Présentation du projet}


Cette section débutera par une vue d'ensemble du contexte métier dans lequel s'inscrit le projet : l'e-commerce. Elle comprendra une définition du terme, une exploration de ses différents modèles économiques, une explication du modèle de l'e-commerce 5c, ainsi qu'une présentation des différents types d'applications e-commerce. Ensuite, elle présentera le client pour lequel le projet sera réalisé : Nespresso, suivi d'une présentation de leur plateforme de commerce électronique. 

\subsection{Contexte métier}
\subsubsection{Définition de l'e-commerce}

L'e-commerce inclut la vente de produits physiques et dématérialisés en ligne, et permet aux entreprises de vendre exclusivement en ligne, offrant une option supplémentaire aux propriétaires d'entreprises qui auraient auparavant dû choisir entre vendre dans un magasin physique ou vendre en gros.

\subsubsection{Modèle 5C}
Le  modèle  5C \cite{5c} définit  le  commerce  électronique  par  cinq  domaines  d'activités dont  les dénominations commencent par la lettre "C" 

\subsubsection{Modèles économiques de l'e-commerce}
Le commerce en ligne peut prendre diverses formes, allant d'une branche numérique d'une grande chaîne de vente au détail ou d'un magasin physique, à une petite entreprise vendant des produits artisanaux depuis son domicile sur des places de marché en ligne.\par

\subsection{Présentation du client}

Le client est une entreprise suisse spécialisée dans la fabrication des machines à café, des capsules de café, et des accessoires tels que des tasses et des porte-capsules.\par 
Son nom est “Nespresso”, c'est un mot-valise constitué à partir de Nestlé et d’Expresso. Elle a été fondée en 1986 et est devenue par la suite une référence mondiale dans l’industrie du café vu la qualité de ses produits.
\par

\begin{figure}[!h]
    \centering
    \includegraphics[width=6cm]{images/nespressoLogo.png}
    \caption{Logo du client Nespresso \cite{nespressowiki}}
    \label{fig:my_label}
\end{figure}

\subsection{Présentation de la plateforme}


Nespresso proposent deux technologies (systèmes de capsules de café) différentes : 
 \begin{enumerate}
     \item La technologie “Original” : la première technologie utilisée par la marque. Elle utilise des capsules en aluminium avec un diamètre de 37 mm et une hauteur de 30 mm. Elle est conçue pour préparer des expressos et des lungos.
     \item La technologie “Vertuo” : une technologie plus récente développée par Nespresso. Elle utilise des capsules en aluminium avec un diamètre plus grand de 62 mm et une hauteur plus courte de 30 mm. Elle est conçue pour préparer une gamme plus large de tailles de tasses de café en utilisant un système de centrifugation breveté.
 \end{enumerate}

\vspace{0.5cm}
Après cette présentation générale du projet, la section suivante abordera la problématique et les objectifs que nous souhaitons atteindre.

\section{Problématique et objectifs}
\subsection{Problématique}

\subsection{Objectifs globaux du projet}
Afin de répondre aux besoins du client en matière d'e-commerce et de scalabilité, SQLI a posé comme objectifs de ce projet : 


\subsection{Objectifs spécifiques}


Ce processus implique une compréhension approfondie de l'architecture existante et de ses limites, ainsi que l'élaboration d'une solution innovante et sur mesure pour répondre aux besoins spécifiques du client en matière de fiabilité, de performance, de sécurité, de réutilisabilité, et d'évolutivité.

\section{Conduite de projet}
 
\subsection{Méthodologie de travail}

\subsubsection{Choix de la méthodologie}

Scrum \cite{scrum} est le cadre de travail agile respecté dans le projet Nespresso. Ainsi, le processus de développement est organisé en sprints, qui sont des itérations courtes et bien définies, au cours desquelles l'équipe se concentre sur la réalisation d'un ensemble de fonctionnalités spécifiques. \par 


\subsubsection{Application sur le projet}

\myparagraph{Software Development Life Cycle (SDLC) 4} 

\myparagraph{Sprint work}

\subsection{Planification du projet}
La planification est une étape importante dans un projet. Elle présente une vision générale sur les étapes de développement et les estimations concernant le délai. Nous avons utilisé le diagramme de GANTT pour présenter le déroulement global du stage de fin d'études 

Le diagramme de Gantt présenté offre une représentation visuelle détaillée de la planification du projet dans la période du stage. 
Nous allons présenter par la suite les résultats des sprints mentionnés.

\section*{Conclusion}
\addcontentsline{toc}{section}{Conclusion}

Dans ce chapitre, nous avons posé les bases du projet en définissant son contexte général, ses objectifs ainsi que le processus de développement adopté pour sa réalisation.

Dans le prochain chapitre, nous procéderons à une étude analytique et conceptuelle approfondie permettant de mieux appréhender la mise en œuvre de la solution.
